\section{Parameter binding}% in the evaluation trace}
\label{sec:parameter-binding}
The preceding discussion treats the core parameter vector
\(\CoreParameters\) as if its coordinates were already known. 
In a typical model-building workflow this identification is not immediate.
A host variable may appear at different levels of the model description and may induce perturbations at many concrete locations in the discrete model. 
A separate registration step is therefore needed to identify which kernel entries are controlled by each scalar
coordinate.

Let
\[
\HostParameters = (\hat y_1,\ldots,\hat y_{n_{\mathrm h}})
\]
denote the scalar host coordinates. 
The model factory is evaluated as
a finite sequence of elementary construction operations defining a computational graph \citep{frostig2018compiling} whose nodes are the intermediate objects created during model construction. 

The set of model objects \(\mathcal{O}\) appearing in this trace is partitioned
into materials, sections, fibers, and elements:
\[
\mathcal{O}
=
\mathcal{O}_{\mathrm{mat}}
\cup
\mathcal{O}_{\mathrm{sec}}
\cup
\mathcal{O}_{\mathrm{fib}}
\cup
\mathcal{O}_{\mathrm{ele}}
\]
The trace induces a
directed dependency relation on \(\mathcal{O}\). 
For example, a fiber depends on the section in which it is created, and an element depends on the section instances referenced in its definition.

\paragraph{Labeling and Binding.}
Each scalar host coordinate is treated together with an identifying
label,
\[
\pi_i = (i,\nu_i,\hat y_i),
\qquad i=1,\ldots,n_{\mathrm h},
\]
where \(i\) is a global differentiation index and \(\nu_i\) is a
symbolic name. The value \(\hat y_i\) is used in the primal
construction, while \((i,\nu_i)\) records the identity of the
coordinate in the differential structure.

% \paragraph{Binding.}
When a labeled coordinate \(\pi_i\) appears in the evaluation trace,
it is bound to an object
\[
a_i \in \mathcal{O},
\]
called its anchor. The anchor specifies the level at which the
coordinate enters the model description. 
A material constant may be anchored at
a material object, a geometric parameter at a section object, and a
fiber area at a fiber object. 

The evaluation trace then propagates this dependence downstream
through the computational graph. 
In this way, one obtains the set of all concrete parameter locations in the kernel that are controlled by
coordinate \(i\).

\paragraph{Kernel addresses.}
Let \(\mathcal{A}\) denote the set of admissible kernel addresses. 
An
address \(\alpha \in \mathcal{A}\) identifies a parameter slot
in the assembled model, for example a material parameter at a given
section point within a given element, or a fiber attribute within a
section occurrence. 
The binding process determines a set-valued map
\[
\mathcal{B} : \{1,\ldots,n_{\mathrm h}\} \to 2^{\mathcal{A}},
\]
where \(\mathcal{B}(i)\) is the set of all kernel addresses reached by
following the dependency graph outward from the anchor \(a_i\) 
and \(2^{\mathcal{A}}\) is the power set of \(\mathcal{A}\) \citep{halmosNaiveSetTheory1998}.


Once the binding map \(\mathcal{B}\) has been resolved, the host
coordinates acquire a concrete realization in the kernel. 
In the
simplest case one identifies
\[
\CoreParameters = (y_1,\ldots,y_{n_{\mathrm c}}),
\qquad
n_{\mathrm c}=n_{\mathrm h},
\qquad
y_i=\hat y_i,
\]
so that the host and kernel use the same scalar coordinates, but the
binding map \(\mathcal{B}\) specifies where each coordinate acts in the
assembled model.

% With this interpretation, the derivative
% \[
% \frac{\partial \ResponseMap_{\ModelFamily}}{\partial y_i}
% \]
% means the sensitivity obtained by perturbing simultaneously every
% kernel address in \(\mathcal{B}(i)\), while holding all remaining
% coordinates fixed. Thus the kernel derivative is not attached to a
% single local slot, but to the full support of a global coordinate
% throughout the evaluation trace.

\paragraph{Forward contraction.}
The forward differential therefore takes the form
\[
\D \ResponseMap_{\ModelFamily}(\CoreParameters)[\dot{\CoreParameters}]
=
\sum_{i=1}^{n_{\mathrm c}}
\frac{\partial \ResponseMap_{\ModelFamily}}{\partial y_i}(\CoreParameters)\,
\dot y_i,
\]
where each partial derivative already accounts for all bound kernel
addresses associated with coordinate \(i\). The host layer determines
the perturbation coefficients \(\dot y_i\), while the kernel layer
evaluates the corresponding directional effect on the response.


\begin{remark}
The preceding description assumes that the evaluation trace produces a
fixed model topology. 
Under this assumption the binding map \(\mathcal{B}\) is locally constant, and the composite response is differentiable in the usual sense. 
If the topology itself changes with
the host parameters, then the map \(\ParameterMap\) is only piecewise
smooth, and the resulting differential structure requires additional care.
\end{remark}